\chapter{Referensi Teknis Python \& PyTorch}
\label{chap:tech_ref}

Sebagai praktisi AI, sintaks bahasa pemrograman adalah alat utama Anda. Bab ini menyediakan referensi teknis padat untuk Python modern dan PyTorch, kerangka kerja Deep Learning paling populer.

\section{Python Modern untuk AI}

\subsection{Type Hinting}
Python 3.10+ mendukung pengetikan statis yang membantu mencegah bug.
\begin{codeblock}[language=Python]
def process_data(items: list[str]) -> dict[str, int]:
    result: dict[str, int] = {}
    for item in items:
        result[item] = len(item)
    return result
\end{codeblock}

\subsection{List \& Dict Comprehensions}
Cara ringkas dan cepat untuk membuat struktur data.
\begin{codeblock}[language=Python]
# List: Kuadrat bilangan genap
numbers = [1, 2, 3, 4, 5]
squares = [x**2 for x in numbers if x % 2 == 0]

# Dict: Map nama ke panjang nama
names = ["Ali", "Budi", "Cici"]
name_lengths = {name: len(name) for name in names}
\end{codeblock}

\subsection{Walrus Operator (:=)}
Melakukan assignment di dalam ekspresi.
\begin{codeblock}[language=Python]
# Tanpa Walrus
data = get_data()
if data:
    process(data)

# Dengan Walrus
if (data := get_data()):
    process(data)
\end{codeblock}

\section{NumPy: Komputasi Matriks}
NumPy adalah fondasi dari semua library AI (PyTorch, TensorFlow, Pandas).

\begin{codeblock}[language=Python]
import numpy as np

# Membuat Array
a = np.array([1, 2, 3])
b = np.ones((3, 3))     # Matriks 3x3 isi 1
c = np.eye(3)           # Matriks identitas

# Operasi Element-wise
x = np.array([1, 2])
y = np.array([3, 4])
print(x * y) # Output: [3, 8]

# Dot Product (Perkalian Matriks)
m1 = np.random.rand(2, 3) # 2 baris, 3 kolom
m2 = np.random.rand(3, 2) # 3 baris, 2 kolom
result = np.dot(m1, m2)   # Hasil: 2x2
\end{codeblock}

\section{PyTorch: Deep Learning Framework}

\subsection{Tensor Basics}
Tensor mirip array NumPy, tapi bisa berjalan di GPU.
\begin{codeblock}[language=Python]
import torch

# Pindah ke GPU jika ada
device = torch.device("cuda" if torch.cuda.is_available() else "cpu")

x = torch.tensor([1.0, 2.0], requires_grad=True).to(device)
y = x * 3 + 2

# Automatic Differentiation (Autograd)
z = y.sum()
z.backward() # Hitung gradien dz/dx
print(x.grad) # Output: [3.0, 3.0]
\end{codeblock}

\subsection{Membangun Neural Network (nn.Module)}
Struktur standar untuk mendefinisikan model.

\begin{codeblock}[language=Python]
import torch.nn as nn
import torch.nn.functional as F

class SimpleNet(nn.Module):
    def __init__(self):
        super().__init__()
        # Layer definisi
        self.fc1 = nn.Linear(784, 128)
        self.fc2 = nn.Linear(128, 10)
        self.dropout = nn.Dropout(0.2)

    def forward(self, x):
        # Aliran data
        x = torch.flatten(x, 1)
        x = F.relu(self.fc1(x))
        x = self.dropout(x)
        x = self.fc2(x)
        return x
\end{codeblock}

\subsection{Training Loop}
Pola standar untuk melatih model.

\begin{codeblock}[language=Python]
model = SimpleNet().to(device)
optimizer = torch.optim.Adam(model.parameters(), lr=1e-3)
criterion = nn.CrossEntropyLoss()

for epoch in range(10):
    for batch, (inputs, targets) in enumerate(dataloader):
        inputs, targets = inputs.to(device), targets.to(device)

        # 1. Forward Pass
        outputs = model(inputs)
        loss = criterion(outputs, targets)

        # 2. Backward Pass & Update
        optimizer.zero_grad() # Reset gradien lama
        loss.backward()       # Hitung gradien baru
        optimizer.step()      # Update bobot

    print(f"Epoch {epoch}, Loss: {loss.item()}")
\end{codeblock}

\subsection{Saving \& Loading}
\begin{codeblock}[language=Python]
# Save (Hanya bobot, lebih aman)
torch.save(model.state_dict(), "model_weights.pth")

# Load
model = SimpleNet() # Inisialisasi arsitektur yang sama
model.load_state_dict(torch.load("model_weights.pth"))
model.eval() # Set mode evaluasi (matikan dropout)
\end{codeblock}

\section{Tips Debugging AI}

\begin{itemize}
    \item **Shape Error:** Paling sering terjadi. Selalu print `tensor.shape` sebelum operasi perkalian matriks. Pastikan `(A, B) x (B, C)`.
    \item **Device Error:** `Expected all tensors to be on the same device`. Pastikan input dan model sama-sama di CPU atau GPU.
    \item **NaN Loss:** Biasanya karena Learning Rate terlalu besar atau pembagian dengan nol. Cek data input untuk nilai tak wajar.
    \item **Overfitting:** Jika `\texttt{training\_loss}` turun tapi `\texttt{validation\_loss}` naik. Solusi: Tambah data, gunakan Dropout, atau perkecil model.
\end{itemize}
